\section{The \texorpdfstring{\(1{:}2\)}{1:2} resonance and quartic near-degeneracy}
\label{sec:resonance}

\subsection{Symmetric primary periodic branch}

Starting from
\[
(0,a,0),
\]
reversibility reduces the periodic-orbit computation to the quarter-period conditions
\[
y(T_{1/4})=0,\qquad
z(T_{1/4})=\frac{\pi}{2}.
\]
Continuation in \(\mu\) produces a smooth primary branch.

\subsection{Numerical localization of the resonance}

The candidate \(1{:}2\) resonance is located by
\[
\tr DP_\mu(0)=-2.
\]
A multiprecision computation gives
\begin{equation}
\boxed{
\mustar
=0.66286900693051012809369086386036365\ldots.
}
\label{eq:mu-star}
\end{equation}
At 120-digit working precision, with an independent high-order Taylor integration of the state and variational equations, the complete event-corrected monodromy satisfies
\[
\|DP_{\mustar}(0)+I\|_F\simeq6.24\times10^{-47},
\]
while
\[
|\det DP_{\mustar}(0)-1|\simeq1.89\times10^{-49}.
\]
The resonance is therefore numerically identified as semisimple.  We summarize this as
\begin{equation}
\boxed{
DP_{\mustar}(0)=-I
}
\label{eq:resonance}
\end{equation}
within the quoted numerical accuracy.  The critical second return
\[
F_0=P_{\mustar}^2
\]
accordingly satisfies \(DF_0(0)=I\) to the same accuracy.

\subsection{Reversibility on the section}

The first reversor induces
\[
r(q,p)=(q,-p).
\]
A complementary reversing involution is generated by
\[
s_\mu=P_\mu r.
\]
Their tangent fixed sets at resonance define two complementary canonical directions. These directions fix the remaining rotational freedom after detuning normalization.

\subsection{Intrinsic detuning}

On the elliptic side define
\[
\delta(\mu)
=
\arccos\left(-\frac{\tr DP_\mu(0)}{2}\right)\ge0.
\]

Let
\[
M_\mu=-P_\mu.
\]
The orientation of \(\delta\) is chosen so that
\[
DM_\mu(0)
=
I+\delta
\begin{pmatrix}
0&-1\\
1&0
\end{pmatrix}
+\bigO(\delta^2).
\]

With
\[
\omega=dQ\wedge dP,\qquad
\iota_{X_H}\omega=dH,
\]
the quadratic generator is
\begin{equation}
\boxed{
H_2=-\delta I,\qquad
I=\frac{Q^2+P^2}{2}.
}
\label{eq:H2}
\end{equation}

\subsection{Reversible linear normalization}

The detuning determines the anisotropic symplectic scale
\[
q=\sqrt{\alpha_{\rm lin}}\,Q,\qquad
p=\frac{P}{\sqrt{\alpha_{\rm lin}}},
\]
with
\[
\alpha_{\rm lin}\simeq0.0946271414514.
\]

The remaining canonical rotations commuting with the normalized linear generator are removed by aligning the \(Q\)- and \(P\)-axes with the tangent fixed sets of the two reversing involutions. The reversible chart is then fixed up to discrete sign changes and interchange of the axes.

If the pure quartic coefficients before the intrinsic scaling are \(a_4q^4+c_4p^4\), the independent balancing scale is
\[
\alpha_{\rm bal}
=
\left(\frac{c_4}{a_4}\right)^{1/4}.
\]
Numerically,
\[
\alpha_{\rm bal}\simeq\alpha_{\rm lin}
\]
to relative order \(10^{-11}\).

\subsection{Odd resonant normal form}

At resonance,
\[
P_\star(w)=-w+\bigO(|w|^2).
\]
We fix the nonlinear gauge by using only the nonresonant canonical generators required to remove the even map terms: a cubic Hamiltonian generator eliminates the quadratic map term and a quintic generator eliminates the quartic map term. No resonant quartic Hamiltonian generator is added. The resulting canonical near-identity transformation puts the map into the odd form
\[
\widetilde P(w)
=
-w+\widetilde P_3(w)+\widetilde P_5(w)+\bigO(|w|^7).
\]

Define
\[
M=-\widetilde P.
\]
Then
\[
M
=
I+M_3+M_5+\bigO(7)
=
\exp(X_H),
\]
with
\[
H=H_4+H_6+\bigO(8).
\]
The extraction is detailed in \cref{app:normal-form}.

\subsection{Quartic interpolating Hamiltonian}

In the fixed reversible chart,
\begin{equation}
\boxed{
H_4
=
A(Q^4+P^4)+BQ^2P^2.
}
\label{eq:H4}
\end{equation}
The currently extracted values are approximately
\[
A\simeq203.6723580,\qquad B\simeq-408.3087120.
\]

Define
\begin{equation}
\boxed{
\rho^2=2-\frac BA.
}
\label{eq:rho-def}
\end{equation}
Then
\[
\rho\simeq2.001182918,
\]
and
\begin{equation}
\boxed{
\sigma=4-\rho^2\simeq-4.7330723\times10^{-3}.
}
\label{eq:sigma-def}
\end{equation}

At \(\rho=2\),
\[
H_4=A(Q^2-P^2)^2.
\]

In polar variables
\[
Q=\sqrt{2I}\cos\theta,\qquad
P=\sqrt{2I}\sin\theta,
\]
one obtains
\[
H_4
=
4AI^2-A\rho^2I^2\sin^22\theta.
\]
Hence
\[
H_4^{\rm ax}=4AI^2,
\]
whereas along an exact diagonal,
\[
H_4^{\rm diag}=A\sigma I^2.
\]

The quartic term is therefore fully active on the axes but strongly suppressed near the diagonals.
